\documentclass{article}
\usepackage[utf8]{inputenc}
\usepackage[T1]{fontenc}
\usepackage{geometry}
\usepackage{graphicx}
\usepackage{float}
\usepackage{hyperref}
\usepackage[english,indonesian]{babel}

\title{Optimasi Pemrosesan Template LaTeX Menggunakan Pendekatan Berbasis AI}
\author{Academic Support Agent}
\date{\today}

\begin{document}
\maketitle

\begin{abstract}

\end{abstract}

\chapter{PENDAHULUAN}


\chapter{TINJAUAN PUSTAKA}


% Original content (commented out):
% \documentclass[a4paper]{article}

\usepackage[pass]{geometry}
\usepackage{CJK,CJKspace,CJKpunct}
\usepackage{amsmath}
\usepackage{graphicx}
\usepackage[colorinlistoftodos]{todonotes}

\begin{CJK*}{UTF8}{zhsong}
\title{論文}
\author{名前}
\end{CJK*}

\pdfmapline{=unisong@Unicode@ <ipam.ttf}
\begin{document}
\begin{CJK*}{UTF8}{zhsong}
\maketitle

\begin{abstract}
あぶすと
\end{abstract}

\section{Introduction}

ここにイントロを書くよ〜
\section{Some \LaTeX{} Examples}
\label{sec:examples}

\subsection{Comments}

サブセクション

\begin{figure}
\centering
\includegraphics[width=0.4\textwidth]{frog.jpg}
\caption{\label{fig:frog}カエルさんだよ〜.}
\end{figure}

\end{CJK*}
\end{document}

\end{document}